% Redesign of the Spring 2026 source deck, with an added Shai et al. example.
% Build from this directory: latexmk -pdf -outdir=build CompMechSlidesPt1.tex
\documentclass[aspectratio=169,12pt,t]{beamer}
\usepackage{iliad-slides}
\graphicspath{{figures/}}
\title{Computational Mechanics}
\subtitle{1. Overview \& Scope}
\institute{ILIAD Intensive}
\date{Autumn 2026}
\author{Xavier Poncini (Simplex)}

\begin{document}
\iliadtitle

% 2. A genuine sequence merits a numbered list.
\begin{frame}{Outline}
  \vspace{5mm}
  \begin{enumerate}
    \setlength{\itemsep}{4mm}
    \item What is Computational Mechanics?
    \item Computational Mechanics \& AI Safety
    \item Overview: where is the program at?
    \item Scope: what will we tackle today?
  \end{enumerate}
\end{frame}

% 3. Keep the two explanations and examples in the same view.
\begin{frame}{What is Computational Mechanics?}
  \begingroup\small
  \keyterm{Statistical Mechanics} makes predictions about the behaviour of
  systems by considering their \keyterm{statistical} properties. E.g.\par\endgroup
  \vspace{-1mm}
  \begin{center}\evidenceimage[17mm]{statmech}\end{center}
  \vspace{-2mm}
  {\small
  \keyterm{Computational Mechanics} (comp-mech) makes predictions about the
  behaviour of systems by considering their \keyterm{computational} properties. E.g.\par}
  \vspace{-1mm}
  \begin{center}\evidenceimage[20mm]{compmech}\end{center}
\end{frame}

% 4. Questions are complete sentences, with no unnecessary hierarchy.
\begin{frame}{Information stored by a system}
  Some questions that are central to comp-mech:
  \vspace{3mm}

  How much of the past does a system store?

  How does the system store this information?

  How does this information affect system behaviour?
  \vspace{4mm}

  We will consider each of these questions for transformers today!
\end{frame}

\begin{frame}{Optimal prediction}
  Not all systems are interesting! Comp-mech places an emphasis on those
  performing \keyterm{optimal prediction}. In this setting the central question becomes:
  \vspace{5mm}

  {\large\color{iliadAccent}What are convergent structures relevant for
  optimal prediction?\par}
  \vspace{5mm}

  Insight on this question greatly constrains our search space when
  considering \keyterm{optimal predictors}.
\end{frame}

\begin{frame}{Instrumental and terminal goals}
  We have strong reason to believe that AI systems will pursue a range of
  \keyterm{instrumental goals} in service of their \keyterm{terminal goal}.
  \begin{center}\evidenceimage[44mm]{instrumental}\end{center}
\end{frame}

\begin{frame}{Optimal prediction as an intervention target}
  Comp-mech uses an \keyterm{instrumental goal} as an intervention target point.
  \begin{center}\evidenceimage[37mm]{intervention}\end{center}
  By considering internal representations consistent with optimal prediction
  the search space is reduced.
\end{frame}

\begin{frame}{Capability scale and optimal prediction}
  Targeting an instrumental goal is \keyterm{scale-free}: it will be relevant
  to current and future models.
  \begin{center}\evidenceimage[23mm]{capability_scale}\end{center}
  {\small
  Which contrasts \keyterm{scale-sensitive} approaches e.g.\ AI control that are
  only applicable up to a fixed capability scale.

  Moreover, as capabilities increase models will better approximate
  \keyterm{optimal predictors}.\par}
\end{frame}

\begin{frame}{The research program}
  Comp-mech applied to AI safety is a young research program
  ($\sim$2 years). There is much to do!

  Progress can be broken down into two orthogonal directions:
  \vspace{2mm}

  For what classes of data generating \keyterm{processes} do we understand
  optimal prediction?

  For what scale of \keyterm{neural networks} have we identified the
  structures of optimal prediction?
\end{frame}

\begin{frame}{Processes and generation resources}
  Fix a (computable) stochastic process and consider:

  What resources are required to generate this process?
  \begin{center}\evidenceimage[44mm]{hierarchy}\end{center}
\end{frame}

\begin{frame}{Processes and optimal prediction}
  For what classes of data generating \keyterm{processes} do we understand
  optimal prediction?
  \begin{center}\evidenceimage[48mm]{hierarchy_green}\end{center}
\end{frame}

% 12. All four quantities share an aligned column and an explicit unit.
\begin{frame}{Neural networks: the scale studied so far}
  For what scale of \keyterm{neural networks} have we identified the structures
  of optimal prediction?
  \vspace{2mm}

  \begin{tabular*}{0.80\linewidth}{@{}l@{\extracolsep{\fill}}r@{}}
    \toprule
    Architecture & Parameters (approximately)\\
    \midrule
    Transformer & $10^{5}$\\
    LSTM        & $10^{5}$\\
    RNN         & $10^{4}$\\
    GRU         & $10^{5}$\\
    \bottomrule
  \end{tabular*}
  \vspace{3mm}

  Very toy compared to current frontier models ($\sim 10^{12}$)!
\end{frame}

% 13. Empirical example following the comparison of network sizes.
\begin{frame}{Example: belief state geometry}
  \setlength{\parskip}{2pt}
  \vspace{-1mm}
  {\centering\evidenceimage[48mm]{shai-belief-state-geometry}\par}
  \vspace{1mm}
  {\small
  After next-token training on HMM sequences, an affine readout approximately
  recovers the geometry of beliefs over hidden states from residual activations.\par}
  \source{Shai et al., \emph{Transformers represent belief state geometry in their
  residual stream}. NeurIPS 2024.
  \href{https://arxiv.org/abs/2405.15943}{arXiv:2405.15943}.}
  \note{A belief state is a probability distribution over the hidden states of
  the data-generating HMM, conditioned on the observed tokens. The lower path
  computes these beliefs by Bayesian updating. The upper path fits an affine
  map from the trained transformer's residual activations to those beliefs.
  The paper finds the predicted geometry, including fractal examples, in the
  studied HMM-trained transformers. Depending on the process, the readout uses
  one layer or concatenated layers. The map is an analysis fitted after training,
  not an extra component of the transformer. The supplied diagram is a schematic,
  not a quantitative plot of reconstruction error. Source: Shai et al., NeurIPS
  2024, Sections 2.3 and 3, https://arxiv.org/abs/2405.15943.}
\end{frame}

\begin{frame}{The research frontier}
  We are working to push the Pareto frontier:
  \begin{center}\evidenceimage[53mm]{pareto}\end{center}
\end{frame}

% 15 (original 14). Preserve diagram and qualification without obscuring data.
\begin{frame}{The research frontier}
  We are working to push the Pareto frontier:
  \vspace{1mm}
  \begin{columns}[T,onlytextwidth]
    \begin{column}{0.49\textwidth}
      \centering\evidenceimage[51mm]{pareto}
    \end{column}
    \begin{column}{0.46\textwidth}
      \vspace{3mm}
      \slidelabel{Wild Speculation}
      As the research program is scale-free one view is that we should
      develop the foundations such that future automated AI Safety
      researchers can pursue these ideas independently.
    \end{column}
  \end{columns}
\end{frame}

\begin{frame}{Scope: what will we tackle today?}
  We will take a slightly narrower view and consider:
  \begin{center}\evidenceimage[53mm]{scope}\end{center}
\end{frame}

\begin{frame}{Plan}
  \vspace{5mm}
  {\small
  \renewcommand{\arraystretch}{1.5}
  \begin{tabularx}{\linewidth}{@{}r@{\hspace{3mm}}X@{\hspace{3mm}}l@{}}
    1 & Predicting the future & Lecture/exercises\\
    2 & Representing the past & Lecture/exercises\\
    3 & Belief geometry in transformers & Reading/Discussion\\
  \end{tabularx}
  \par}
\end{frame}

\begin{frame}[plain,c]
  \leavevmode
  {\usebeamerfont{title}Questions?\par}
\end{frame}
\end{document}
